\section{The central value and the exact formula}\label{sec:center}

The unit comparison now identifies the central valuation of the
strict determinant with the logarithmic expression of
\eqref{eq:source-12}.  We compute that determinant in the full Tate
lattice.  This is where the point index, torsion, and every local
component number enter the equality.

\subsection{The central determinant volume}

Keep $p>2$ fixed.  Put $K=\Q(\sqrt D)$ as in
Lemma~\ref{classical:auxiliary-fields}, and let
$P_K\in E(K)$ be the single Hilbert-class trace used in
\eqref{eq:source-12}.  The simple zero of $L(E/K,s)$ and
Theorem~\ref{classical:gross-zagier} make $P_K$ nontorsion.
The group $E(K)$ has rank one and $\Sha(E/K)$ is finite: the two
rational factors have analytic ranks adding to one, and the
quadratic restriction-of-scalars isogeny transfers rank and
finiteness.  Define
\[
 n_P=v_p\bigl([E(K)/E(K)_{\rm tors}:\Z\overline P_K]\bigr),
 \quad \tau_g=v_p(\#E(K)_{\rm tors}),
 \quad s_K=v_p(\#\Sha(E/K)).
\]
For a rational finite prime $q$ write
$\tau_q=v_p(\#E(\Q_q)[p^\infty])$, and write
$\log_{\omega_E}E(\Q_p)=p^l\Z_p$.
All primes in the fixed support $S^\circ$ of the $E/K$ comparison
are split in $K$; thus their two local quantities equal the
corresponding rational quantities.  The modular differential
factor is the actual rational number $c_E$ in
$\phi^*\omega_E=c_E f\,dq/q$.

\begin{proposition}\label{center:exact-factor}
Let $C_w^{s=1}(0)$ be the actual central strict/full complex after
all artificial variables have been removed.  Its differential
determinant satisfies
\begin{equation}
 \begin{split}
 -v_p(L_w^{s=1}(0))
 &=d(C_w^{s=1}(0),1)\\
 &=2n_P+2\tau_g-s_K
   -2\bigl(v_p(\log_{\omega_E}P_K)-l+\tau_p\bigr)
   -2\sum_{q\in S^\circ}\tau_q .
 \end{split}
 \tag{22}\label{eq:source-22}
\end{equation}
Consequently
\begin{equation}
 2(n_P+\tau_g)=s_K+2v_p(c_E)+2\sum_q v_p(c_q(E)),
 \tag{23}\label{eq:source-23}
\end{equation}
where the last sum ranges over all rational finite primes.
\end{proposition}

\begin{proof}
After removal of the artificial variables, the marked arithmetic
models compute the fixed-support complex over $K$ itself by
Lemma~\ref{models:finite-models}.  Central specialization is
derived specialization of its finite free model.  In particular
it retains the finite local Kummer modules and their invariant
torsion.  The strict/full complex is rationally acyclic: the
global compact Selmer space is the line generated by $P_K$,
and its localization at each $p$-place is nonzero, since the
kernel of the elliptic logarithm on local points is torsion.
Exact complementary duality, or
Lemma~\ref{models:strict-finite}, then gives the acyclicity.
The inverse-determinant convention consequently gives the first
line of \eqref{eq:source-22}.

Begin with the all-Kummer complex.  By
Lemma~\ref{models:kummer-volume}, its class-functional tensor has
valuation
\[
 d(C_{\rm Kum},(P_K,P_K^\vee))=2n_P+2\tau_g-s_K.
\]
Here the torsion lengths in degrees one, two, three are
$\tau_g,s_K,\tau_g$, and both occurrences of the point contribute
its index $n_P$.  Thus this is a calculation on the full lattice.

First compare the paired $p$-adic Kummer lines with strict at $w$
and full at $\bar w$, keeping the Kummer conditions at every
non-$p$ place.  Denote this intermediate complex by
$C_w^{\rm Kum-away}$; it differs from $C_w^{s=1}(0)$ only at the
non-$p$ support places.  The unchanged conditions are thus still exact
complements, as required by Lemma~\ref{models:strict-finite}.
Choose a free local Kummer generator $e$
with logarithm generating $p^l\Z_p$.  Localization of $P_K$ has
free coordinate of valuation
$v_p(\log_{\omega_E}P_K)-l$.  The local Kummer complex also
contains degree-one torsion of length $\tau_p$; its volume in
this rational basis is therefore $\tau_p$.  The two paired
localization triangles in
\eqref{models:strict-finite-volume} subtract twice the sum of
these two quantities, giving
\[
 d(C_w^{\rm Kum-away},1)=2n_P+2\tau_g-s_K
       -2\bigl(v_p(\log_{\omega_E}P_K)-l+\tau_p\bigr).
\]

Now relax the conditions at the non-$p$ support places.  At
such a place compact $H^1$ is already the completed point module,
while the quotient of full cochains by the Kummer condition has
degree-two torsion of length $\tau_q$, by local duality.
Its inverse-determinant valuation is $-\tau_q$.  Each rational
$q\in S^\circ$ gives two places of $K$, hence a total change
$-2\sum_{q\in S^\circ}\tau_q$.  These last comparisons use
determinant multiplication in localization triangles and do not
require full local conditions to be self-complementary.
Combining the two changes proves \eqref{eq:source-22}.
In particular the torsion term $\tau_p$
cannot be discarded by replacing the point module by its free
quotient.

By Theorem~\ref{unit:unit-comparison} the analytic series and
strict determinant on the true line differ by an integral unit.
Their values at $t=0$ are therefore of the same valuation.
Equation~\eqref{eq:source-12} gives
\[
 v_p(L_w^{s=1}(0))
 =-2v_p(c_E)+2v_p(\log_{\omega_E}P_K)
   +2v_p(P_p(1))+2\sum_{q\in S^\circ}v_p(P_q(1)).
\]
Every argument of $v_p$ in this display is nonzero.  In particular $P_q(1)$ is a
nonzero inverse local Euler factor, and the point logarithm is
nonzero as observed above.

Set $h=v_p(\log_{\omega_E}P_K)$, $e_q=v_p(P_q(1))$, and
$c=v_p(c_E)$.  Equating this last expression with the negative
of \eqref{eq:source-22}, then cancelling $2h$, gives
\[
 2n_P+2\tau_g
 =s_K+2c+2(\tau_p-l-e_p)
       +2\sum_{q\in S^\circ}(\tau_q-e_q).
\]
The Haar identities of Lemma~\ref{models:haar} identify the
parentheses with $v_p(c_p(E))$ and $v_p(c_q(E))$, respectively.
All bad primes are in the fixed support, and every omitted good
prime has component number one.  This proves
\eqref{eq:source-23} with the sum over all finite primes.
\end{proof}

\begin{corollary}\label{center:pair-discrepancy}
For every elliptic curve of analytic rank zero or one, including
CM curves, and for the imaginary twist in
Lemma~\ref{classical:auxiliary-fields},
\begin{equation}
                         X(E)+X(E^D)=0.
                         \tag{24}\label{eq:source-24}
\end{equation}
\end{corollary}

\begin{proof}
The full group index in
Proposition~\ref{classical:discrepancy-identity} has valuation
\[
 v_p([E(K):\Z P_K])=n_P+\tau_g.
\]
Substitute \eqref{eq:source-23} into the resulting valuation
identity \eqref{eq:classical-discrepancy}.  Its right-hand side
is zero, proving \eqref{eq:source-24}.  That arithmetic-volume
identity already uses the total real periods, full-lattice
regulator, and height diagonal $H$.  The field-degree factor
over $K$ is included there; none is appended here.
\end{proof}

\subsection{CM curves and assembly}

The single-curve inequality applies to non-CM curves.  The
rank-zero CM case instead has the following independent input,
which we state at its exact scope.

\begin{lemma}\label{center:cm-zero}
If $E/\Q$ has complex multiplication and $L(E,1)\ne0$, the exact
BSD formula holds for $E$ with the normalizations of
Section~\ref{sec:intro}.
\end{lemma}

\begin{proof}
Let $M$ be the CM field.  Theorem 1.1 and Corollary 1 of
Burungale--Flach~\cite{BurungaleFlach} give exact BSD for an
elliptic curve $A/F$ with CM by the maximal order of $M$, provided
$F(A_{\rm tors})/M$ is abelian and the associated Hecke
$L$-value at one is nonzero.  We verify these hypotheses after an
isogeny over $M$.

All CM endomorphisms of $E_M$ are defined over $M$.  If the
endomorphism order is not maximal, at each prime $\ell$ replace
$T_\ell E$ by its span under $\mathcal O_M\otimes\Z_\ell$ in
$V_\ell E$.  These are Galois-stable lattices, equal to the
original lattices except at the finitely many primes dividing
the conductor of the order.  Their finite quotients define a
finite Galois-stable subgroup of $E$, and quotienting gives an
elliptic curve $A/M$ isogenous to $E_M$ with maximal-order CM.
Its Tate modules have rank one over the corresponding local CM
orders.  The Galois actions commute with those orders, so their
images, and hence the full torsion extension over $M$, are
abelian.

Choose a nonzero trace-zero CM endomorphism after multiplying
by an integer to make it integral.  Galois conjugation outside
$M$ sends it to its negative.  The twisting descent criterion
therefore gives a $\Q$-isogeny from the quadratic twist $E^M$ to
$E$.  In particular
\[
 L(E/M,s)=L(E/\Q,s)L(E^M/\Q,s)=L(E/\Q,s)^2.
\]
This has nonzero value at one.  Isogeny does not change that
$L$-function, and its factorization into the Hecke character
and conjugate Hecke character $L$-functions shows that the
nonvanishing hypothesis of Burungale--Flach holds for $A$.
Their theorem gives exact BSD over $M$, and isogeny invariance
transfers it to $E_M$.

Finally restriction of scalars is compatible with the BSD
arithmetic volume, and
\[
 \operatorname{Res}_{M/\Q}E_M\ \sim_{\Q}\ E\times E^M
 \ \sim_{\Q}\ E^2.
\]
Thus the normalized BSD discrepancy for $E$ has square one.
It is positive: analytic rank zero gives rank zero and finite
$\Sha$, so the two-primary result supplies positivity
and rationality.  It is therefore one.  This is the descent
argument of Burungale--Flach, Corollary 2, using the
restriction-of-scalars and isogeny compatibilities of
\cite{Cassels,MilneBSD}; the preceding lattice argument supplies
the additional nonmaximal-order case.
\end{proof}

\begin{proposition}\label{center:odd-valuations}
For every elliptic curve $E/\Q$ of analytic rank zero or one and
every odd prime $p$,
\[
                  v_p(Q_E)=v_p(\#\Sha(E/\Q)).
\]
\end{proposition}

\begin{proof}
For a non-CM curve, the single-curve inequality
\eqref{eq:source-8} gives $X(E)\geq0$.  Its twist $E^D$ is
also non-CM and has the complementary analytic rank, so the
same inequality gives $X(E^D)\geq0$.  Their sum is zero by
\eqref{eq:source-24}; hence $X(E)=0$.

For a CM curve of analytic rank zero use
Lemma~\ref{center:cm-zero}.  For a CM curve of analytic rank one,
the companion $E^D$ is CM of analytic rank zero.  That lemma and
\eqref{eq:source-24} again give $X(E)=0$.  The construction was
made for an arbitrary odd $p$, including $p=3$, so the displayed
equality holds at every odd prime.
\end{proof}

\begin{proof}[Proof of Theorem~\ref{thm:main}]
Apply the low-corank converse
(Theorem~\ref{classical:low-corank}) at the prime $q$ in the
hypothesis.  It gives
$a(E)=r(E)=s_q(E)\in\{0,1\}$ and finiteness of the whole
$\Sha(E/\Q)$.  The Kummer exact sequence then gives
$s_p(E)=r(E)$ for every prime $p$, in particular $s_2(E)\leq1$.
The two-primary formula
(Theorem~\ref{classical:two-primary}) gives positivity and
rationality of $Q_E$ and its correct $2$-adic valuation.
Proposition~\ref{center:odd-valuations} supplies its correct
valuation at every odd prime, independently of $q$.
Consequently the positive rational number
$Q_E/\#\Sha(E/\Q)$ has zero valuation at every rational prime.
Its prime factorization is empty, so it equals one.  Substituting
the definition of $Q_E$ is exactly the asserted real formula,
with the total real period and full Mordell--Weil regulator.
\end{proof}
